2.4. \underline{La categoría $\operatorname{BIEXT}(P,Q,G)$ para
$P,Q,G$ variables}.

Afirmo que \underline{el funtor canónico (2.3.4) es fibrado} y que
\underline{el funtor canónico (2.3.5) es cofibrado}. La primera
afirmación significa, en particular, que si $E$ es una biextensión de
$(P,Q)$ por $G$, y si
\[
 u:P'\longrightarrow P,\qquad v:Q'\longrightarrow Q
\]
son morfismos de grupos abelianos, existe una biextensión $E'$ de
$(P',Q')$ por $G$ y un morfismo de biextensiones
\[
 f:E'\longrightarrow E
\]
compatible con $u,v,\mathrm{id}_G$, que posee la propiedad universal
habitual respecto de los morfismos de biextensiones hacia $E$
compatibles con $u,v,\mathrm{id}_G$.

La segunda afirmación se explicita de manera análoga. En particular,
si $E$ es una biextensión de \sisi{$(P,G)$}{$(P,Q)$} por $G$, y si
\[
 w:G\longrightarrow G'
\]
es un morfismo de grupos abelianos,

%%%%%%%%%%  scan idx 175  |  folio 164  %%%%%%%%%%

existe una biextensión $E'$ de $(P,Q)$ por $G'$ y un morfismo de
biextensiones
\[
 g:E\longrightarrow E'
\]
compatible con $\mathrm{id}_P,\mathrm{id}_Q,w$, que posee la propiedad
universal habitual respecto de los morfismos de biextensiones desde
$E$ compatibles con $\mathrm{id}_P,\mathrm{id}_Q,w$. Dejamos al lector
la explicitación de las dos construcciones cuya existencia acaba de
afirmarse; se obtienen fácilmente parafraseando las construcciones
análogas bien conocidas para extensiones de grupos ordinarios.

Así, para tres grupos abelianos $P,Q,G$ de $\underline T$, escribimos
\begin{equation}\tag{2.4.1}
 \operatorname{BIEXT}(P,Q;G)
\end{equation}
para la categoría de biextensiones de
\sisi{$(P;Q)$}{$(P,Q)$} por $G$, definida como la categoría fibra de
(2.3.3) sobre $(P,Q,G)$ (de modo que sus morfismos son los morfismos
de biextensiones compatibles con los morfismos identidad de $P,Q,G$).
Entonces \underline{esta categoría depende covariantemente de $G$ y
contravariantemente de $P,Q$}. Además, con estas varianzas,
\underline{la categoría $\operatorname{BIEXT}(P,Q;G)$ depende
aditivamente de cada uno de los tres argumentos $P,Q,G$}. Para el
argumento $G$, por ejemplo, esto significa precisamente:

\begin{enumerate}[label=\alph*)]
\item $\operatorname{BIEXT}(P,Q;0)$ es equivalente a la categoría de
un solo objeto;
\item para dos grupos abelianos $G,G'$ de $\underline T$, el funtor
\begin{equation}\tag{2.4.2}
 \sisi{\operatorname{BIEXT}(P,G;G\times G')}%
       {\operatorname{BIEXT}(P,Q;G\times G')}
 \longrightarrow
 \operatorname{BIEXT}(P,Q;G)\times
 \operatorname{BIEXT}(P,Q;G')
\end{equation}
deducido de las dos proyecciones $G\times G'\longrightarrow G$ y
$G\times G'\longrightarrow G'$ es una
\underline{equivalencia de categorías}.
\end{enumerate}
En el lenguaje de [4], se dice también que, para $P,Q$ fijos, la
categoría de biextensiones de $(P,Q)$ por grupos variables $G$ es,

%%%%%%%%%%  scan idx 176  |  folio 165  %%%%%%%%%%

mediante el funtor $E\mapsto G$ inducido por (2.3.5), una
\underline{categoría cofibrada aditiva} sobre
$\operatorname{GRAB}(\underline T)$.

\underline{Observación 2.4.3.} En la terminología de \emph{loc.\ cit.},
la categoría cofibrada precedente no solo es aditiva, sino incluso
\underline{exacta por la izquierda} [4, 2.2]. No verificaremos aquí
esta propiedad, pues reaparecerá más adelante en 3.6 como consecuencia
de hechos más precisos. Señalemos únicamente que la teoría general de
\emph{loc.\ cit.} permitiría describir la estructura de esta categoría
cofibrada, salvo equivalencia de categorías fibradas, mediante un
complejo de cadenas $L_\bullet$ en
$\operatorname{GRAB}(\underline T)$ que solo depende de $P$ y $Q$.
Más adelante daremos también una construcción directa de este complejo
como
\[
 P\overset{\mathbb L}{\otimes}Q,
\]
el producto tensorial tomado en una categoría derivada.
